% Language: LaTeX
% =======================  CON BARRA DE TITULO
%   Estandar (1 argumento {texto barra}):
%   codeBlockC         barra oscura, syntax C
%   codeBlockSV        barra oscura, syntax SystemVerilog
%   codeBlockBash      barra oscura, syntax bash
% =======================  Esteticas ([estilo opcional]{texto barra}; por defecto C):
%   codeBlockLat       franja azul lateral + titulo azul (callout)
% # codeBlockTab       pestana oscura arriba-izquierda
%   codeBlockIeee      barra azul IEEE
%   codeBlockTermROG   terminal oscuro + circulos macOS (rojo/naranja/verde)
% # codeBlockTermGGG   terminal oscuro + tres circulos verdes
% * codeBlockTermLinux terminal oscuro + boton [+] de GNOME/Linux
%     -> estilos fondo CLARO  (Clat/Ctab/Cieee):       caja_c | caja_sv | caja_bash
%     -> estilos fondo OSCURO (Cterm/Cterm2/Ctermlnx): caja_c_dark | caja_sv_dark | caja_bash_dark
% =======================  SIN BARRA
%   codeBlockMinC      marco fino sobre blanco; [estilo opcional] (def. caja_c)
%   codeBlockMinCcap      marco sobre fondo gris, syntax C
%   codeBlockMinSVcap     marco sobre fondo gris, syntax SystemVerilog
%   codeBlockMinBashcap   marco sobre fondo gris, syntax bash

% ===================================================================================================== Ejemplo generico
% --- codeBlockTermROG ---  (terminal; aqui en bash con [caja_bash_dark])
% \begin{codigo}
% \begin{TermROGcode}[caja_LENGUAJE_dark]{TITULO}
% CODIGO
% \end{TermROGcode}
% \caption{Ejemplo con codeBlockTermROG.}
% \label{lst:ej-term}
% \end{codigo}
% =====================================================================================================
% * -> presenta problemas de render
% # -> en uso

% =====================================================================================================
%  listings.tex  --  Estilos de codigo fuente
%  Requiere: xcolor, listings, tcolorbox, caption, newfloat (paquetes.tex)
%  Define 'ieeeblue' (compartido); entornos.tex tambien lo usa.
% =====================================================================================================

% --- Colores generales (compartidos por varios formatos) ---
\definecolor{ieeeblue}{HTML}{1F4E79}     % keyword azul (compartido con notaieee)
\definecolor{comentario}{HTML}{5B8A5B}   % comentarios verdes
\definecolor{barra}{HTML}{3A3F44}        % barra de titulo oscura
\definecolor{fondocodigo}{HTML}{F7F7F7}  % cuerpo del codigo
\definecolor{marco}{HTML}{C9CCD1}        % borde sutil
\definecolor{coltitulo}{HTML}{F4F4F4}    % texto de la barra


% =====================================================================================================
%  FLOAT "codigo"  (estilo figure/table; UN solo contador e indice)
%  - \caption dentro del entorno numera "Codigo X.Y."
%  - \label dentro del entorno SI lo detecta el autocompletado
%  - Indice:  \listofcodigo
%  within=section -> "seccion.n"; cambia a 'chapter' en report/book.
% =====================================================================================================
\DeclareFloatingEnvironment[
  fileext=locodigo, listname={Índice de Códigos},
  name=Código, placement=H, within=section]{codigo}
\captionsetup[codigo]{labelsep=period, font=small, format=hang}


% =====================================================================================================
%  FORMATO ANTERIOR  (independiente; numeracion nativa de 'listings')
%  Uso:  \begin{lstlisting}[style=codigo, language=C, caption=..., label=...]
% =====================================================================================================
\lstdefinestyle{codigo}{
  basicstyle=\ttfamily\footnotesize, keywordstyle=\color{ieeeblue}\bfseries,
  commentstyle=\color{gray}\itshape, stringstyle=\color{teal!70!black},
  numbers=left, numberstyle=\tiny\color{gray!70}, numbersep=7pt,
  frame=single, rulecolor=\color{gray!40}, backgroundcolor=\color{gray!4},
  breaklines=true, showstringspaces=false, tabsize=2, captionpos=b,
  xleftmargin=10pt, framexleftmargin=10pt}
\lstdefinestyle{codigo_sv}{style=codigo, language=Verilog,
  morekeywords={
    import,byte,unsigned,logic,bit,void,automatic,function
    }
    }


% =====================================================================================================
%  ESTILOS DE LISTING (sintaxis)  --  compartidos por todas las cajas nuevas
%  Fondo CLARO:  caja_c | caja_sv | caja_bash
%  Fondo OSCURO: caja_c_dark | caja_sv_dark | caja_bash_dark  (terminal)
% =====================================================================================================
\lstdefinestyle{caja_base}{
  basicstyle=\ttfamily\footnotesize,
  keywordstyle=\color{ieeeblue}\bfseries,
  commentstyle=\color{comentario}\itshape,
  stringstyle=\color{teal!70!black},
  numbers=none, frame=none,
  breaklines=true, breakautoindent=false, breakindent=0pt,
  showstringspaces=false, tabsize=2,
  columns=fullflexible, keepspaces=true,
}
\lstdefinestyle{caja_c}{style=caja_base, language=C}
\lstdefinestyle{caja_sv}{style=caja_base, language=Verilog,
  morekeywords={
    import,byte,unsigned,logic,bit,void,automatic,function
    }
    }
\lstdefinestyle{caja_bash}{style=caja_base, language=Bash,
  morecomment=[l]{\#},
  morekeywords={
    if,then,else,elif,fi,for,while,until,do,done,case,esac,in,select,function,break,continue,return,exit,local,export,readonly,unset,source,eval,exec,trap,shift,test,true,false,time,echo,printf,read,cat,tac,head,tail,less,more,grep,egrep,fgrep,sed,awk,cut,paste,sort,uniq,find,xargs,which,whereis,ls,cp,mv,rm,mkdir,rmdir,chmod,chown,ln,touch,pwd,cd,pushd,popd,dirs,tar,gzip,gunzip,zip,unzip,diff,cmp,patch,make,cmake,ninja,gcc,g++,clang,clang++,vcs,iverilog,verilator,vsim,xrun,python,python3,pip,pip3,perl,ruby,node,npm,yarn,git,pull,push,clone,commit,checkout,merge,rebase,fetch,status,branch,tag,curl,wget,ssh,scp,rsync,sudo,su,env,envsubst,sleep,wait,kill,killall,basename,dirname,realpath,seq,expr,yes,eq,ne,lt,le,gt,ge,script
},
  morestring=[b]", morestring=[b]'}


% =====================================================================================================
%  CAJAS DE DISENO  (solo dibujan la caja; el caption/label van en el float)
%  CON barra -> argumentos:  [estilo de listing]{texto de la barra}
%               el estilo es OPCIONAL (por defecto C); ver tabla de estilos arriba
% =====================================================================================================

% --- Estandar: barra oscura -> codeBlockC | codeBlockSV | codeBlockBash ---
\tcbset{estilocaja/.style={
  enhanced, breakable, listing only,
  colback=fondocodigo, colbacktitle=barra, coltitle=coltitulo,
  colframe=marco, boxrule=0.6pt,
  arc=3pt, outer arc=3pt, titlerule=0pt,
  toptitle=2.5pt, bottomtitle=2.5pt, lefttitle=10pt,
  top=5pt, bottom=5pt, left=8pt, right=8pt, boxsep=0pt,
  fonttitle=\normalfont\small, coltext=black}}
\newtcblisting{Ccode}[1]{estilocaja, listing style=caja_c,    title={#1}}
\newtcblisting{SVcode}[1]{estilocaja, listing style=caja_sv,  title={#1}}
\newtcblisting{Bashcode}[1]{estilocaja, listing style=caja_bash, title={#1}}

% --- V1: Barra lateral / callout -> codeBlockLat ([estilo]{barra}) ---
\tcbset{caja_lateral/.style={
  enhanced, breakable, listing only,
  colback=fondocodigo, coltext=black,
  colframe=fondocodigo, boxrule=0pt, sharp corners,
  borderline west={3pt}{0pt}{ieeeblue},
  colbacktitle=fondocodigo, coltitle=ieeeblue, titlerule=0pt,
  toptitle=4pt, bottomtitle=0pt, lefttitle=12pt,
  top=2pt, bottom=4pt, left=12pt, right=8pt, boxsep=0pt,
  fonttitle=\normalfont\small\bfseries}}
\newtcblisting{Latcode}[2][caja_c]{caja_lateral, listing style=#1, title={#2}}

% --- V2: Terminales oscuros -> codeBlockTermROG | codeBlockTermGGG | codeBlockTermLinux ---
%   Comparten estilos de listing oscuros: caja_c_dark (def.) | caja_sv_dark | caja_bash_dark
%   y la base 'caja_terminal_base'. Difieren SOLO en los botones de la barra (overlay).
\definecolor{termbg}{HTML}{282C34}
\definecolor{termfg}{HTML}{ABB2BF}
\definecolor{termkw}{HTML}{61AFEF}
\definecolor{termcm}{HTML}{7F848E}
\definecolor{termst}{HTML}{98C379}
\lstdefinestyle{caja_c_dark}{
  basicstyle=\ttfamily\footnotesize\color{termfg},
  keywordstyle=\color{termkw}\bfseries, commentstyle=\color{termcm}\itshape,
  stringstyle=\color{termst}, identifierstyle=\color{termfg},
  language=C, numbers=none, frame=none,
  breaklines=true, showstringspaces=false, tabsize=2,
  columns=fullflexible, keepspaces=true}
\lstdefinestyle{caja_sv_dark}{style=caja_c_dark, language=Verilog,
  morekeywords={
    import,byte,unsigned,logic,bit,void,automatic,function
    }
    }
\lstdefinestyle{caja_bash_dark}{style=caja_c_dark, language=bash,
  morecomment=[l]{\#},
  morekeywords={
    if,then,else,elif,fi,for,while,until,do,done,case,esac,in,select,function,break,continue,return,exit,local,export,readonly,unset,source,eval,exec,trap,shift,test,true,false,time,echo,printf,read,cat,tac,head,tail,less,more,grep,egrep,fgrep,sed,awk,cut,paste,sort,uniq,find,xargs,which,whereis,ls,cp,mv,rm,mkdir,rmdir,chmod,chown,ln,touch,pwd,cd,pushd,popd,dirs,tar,gzip,gunzip,zip,unzip,diff,cmp,patch,make,cmake,ninja,gcc,g++,clang,clang++,vcs,iverilog,verilator,vsim,xrun,python,python3,pip,pip3,perl,ruby,node,npm,yarn,git,pull,push,clone,commit,checkout,merge,rebase,fetch,status,branch,tag,curl,wget,ssh,scp,rsync,sudo,su,env,envsubst,sleep,wait,kill,killall,basename,dirname,realpath,seq,expr,yes,eq,ne,lt,le,gt,ge,script
    if,then,else,elif,fi,for,while,do,done,echo,rm,cd,git,make,source,exit,code,
    % otros
    testfloat_gen, testfloat_ver, level, seed
    }
  }
% Apariencia comun. coltext/colupper=termfg EVITA texto negro: los identificadores,
% numeros y signos de puntuacion NO toman el color de basicstyle, sino el del entorno.
\tcbset{caja_terminal_base/.style={
  enhanced, unbreakable, listing only,
  colback=termbg, colframe=termbg, boxrule=0pt,
  before upper={\color{termfg}},   % <-- en lugar de: coltext=termfg, colupper=termfg
  colbacktitle=black!88, coltitle=termfg,
  arc=5pt, outer arc=5pt, titlerule=0pt,
  toptitle=3pt, bottomtitle=3pt, halign title=center,
  top=5pt, bottom=5pt, left=10pt, right=8pt, boxsep=0pt,
  fonttitle=\normalfont\footnotesize\ttfamily}}

% V2a: macOS (circulos rojo/naranja/verde) -> codeBlockTermROG
\tcbset{caja_terminal/.style={caja_terminal_base,
  overlay unbroken and first={
    \fill[red!75!black]    ([xshift=12pt,yshift=-9pt]frame.north west) circle (2.3pt);
    \fill[orange!80!black] ([xshift=22pt,yshift=-9pt]frame.north west) circle (2.3pt);
    \fill[green!65!black]  ([xshift=32pt,yshift=-9pt]frame.north west) circle (2.3pt);}}}
\newtcblisting{TermROGcode}[2][caja_c_dark]{caja_terminal, listing style=#1, title={#2}}
% V2a-2: tres circulos verdes -> codeBlockTermGGG
\tcbset{caja_terminal_verde/.style={caja_terminal_base,
  overlay unbroken and first={
    \fill[green!65!black] ([xshift=12pt,yshift=-9pt]frame.north west) circle (2.3pt);
    \fill[green!65!black] ([xshift=22pt,yshift=-9pt]frame.north west) circle (2.3pt);
    \fill[green!65!black] ([xshift=32pt,yshift=-9pt]frame.north west) circle (2.3pt);}}}
\newtcblisting{TermGGGcode}[2][caja_c_dark]{caja_terminal_verde, listing style=#1, title={#2}}

% --- V2b: Terminal oscuro estilo GNOME/Linux -> codeBlockTermLinux ([estilo]{barra}) ---
%   Igual que codeBlockTermROG pero con el boton [+] de Linux en vez de los circulos.
\tcbset{caja_terminal_lnx/.style={caja_terminal_base,
  overlay unbroken and first={
    \fill[green!55!black, rounded corners=1.5pt]
      ([xshift=10pt,yshift=-5pt]frame.north west) rectangle ([xshift=23pt,yshift=-13pt]frame.north west);
    \draw[green!30!black, rounded corners=1.5pt, line width=0.3pt]
      ([xshift=10pt,yshift=-5pt]frame.north west) rectangle ([xshift=23pt,yshift=-13pt]frame.north west);
    \draw[white, line width=0.7pt, line cap=round]
      ([xshift=16.5pt,yshift=-7pt]frame.north west) -- ([xshift=16.5pt,yshift=-11pt]frame.north west);
    \draw[white, line width=0.7pt, line cap=round]
      ([xshift=14.5pt,yshift=-9pt]frame.north west) -- ([xshift=18.5pt,yshift=-9pt]frame.north west);}}}
\newtcblisting{TermLinuxcode}[2][caja_c_dark]{caja_terminal_lnx, listing style=#1, title={#2}}

% --- V3: Pestana -> codeBlockTab ([estilo]{barra}) ---
\tcbset{caja_tab/.style={
  enhanced, unbreakable, listing only,
  colback=fondocodigo, coltext=black, colframe=marco, boxrule=0.6pt,
  arc=3pt, outer arc=3pt,
  top=14pt, bottom=5pt, left=8pt, right=8pt, boxsep=0pt,
  fonttitle=\normalfont\footnotesize\bfseries, coltitle=white,
  attach boxed title to top left={xshift=8pt, yshift=-7pt},
  boxed title style={colback=barra, colframe=barra, arc=2pt, boxrule=0pt,
                     top=1.5pt, bottom=1.5pt, left=5pt, right=5pt}}}
\newtcblisting{Tabcode}[2][caja_c]{caja_tab, listing style=#1, title={#2}}

% --- V4: IEEE azul -> codeBlockIeee ([estilo]{barra}) ---
\tcbset{caja_ieee/.style={estilocaja, colbacktitle=ieeeblue, colframe=ieeeblue!55}}
\newtcblisting{Ieeecode}[2][caja_c]{caja_ieee, listing style=#1, title={#2}}


% =====================================================================================================
%  ESTILOS SIN BARRA SUPERIOR
% =====================================================================================================

% --- V5: Minimalista -> codeBlockMinC (SIN barra; [estilo] opcional) ---
\tcbset{caja_min/.style={
  enhanced, breakable, listing only,
  colback=white, coltext=black, colframe=gray!50, boxrule=0.5pt,
  arc=2pt, outer arc=2pt,
  top=5pt, bottom=5pt, left=8pt, right=8pt, boxsep=0pt}}
\newtcblisting{MinCcode}[1][caja_c]{caja_min, listing style=#1}

% --- Sin decoracion -> codeBlockMinCcap | codeBlockMinSVcap | codeBlockMinBashcap (sin args) ---
\tcbset{cajacap/.style={
  enhanced, breakable, listing only,
  colback=fondocodigo, colframe=marco, boxrule=0.6pt,
  arc=3pt, outer arc=3pt,
  top=5pt, bottom=5pt, left=8pt, right=8pt, boxsep=0pt}}
\newtcblisting{MinCcapcode}{cajacap, listing style=caja_c}
\newtcblisting{MinSVcapcode}{cajacap, listing style=caja_sv}
\newtcblisting{MinBashcapcode}{cajacap, listing style=caja_bash}